%!TEX encoding = UTF-8 Unicode
%!TEX root = reporte-teoria-del-estado-y-constitucion-Actividad-8.tex
% arara: xelatex: { interaction: nonstopmode, synctex: true }
% arara: biber
% arara: xelatex: { interaction: nonstopmode, synctex: true }
% arara: xelatex: { interaction: nonstopmode, synctex: true }

\documentclass[12pt, a4paper, oneside]{article}
\usepackage[aulatex]{template-unadm}
\usepackage{template-unadm-ld}
\usepackage{template-unadm-lex}

% Datos canónicos del estudiante
\def\documentauthor{Mart\'in Jonathan de la Cruz Mu\~noz}
\author{Mart\'in Jonathan de la Cruz Mu\~noz}
\def\documentauthorcanonico{Martín Jonathan de la Cruz Muñoz}  % UTF-8 para validación
\def\documentsubject{Licenciatura en Derecho}
\def\documenttitle{Foro de razón de la existencia del Estado}
\def\documenttitlehf{Foro de razón de la existencia del Estado}
\def\documentsubtitle{Reflexiones sobre los fines del Estado en la era de la globalización}
\def\documentdate{Semana 8, 2026}
\def\documentperiodo{2026-2 B2}
\def\documentcareer{Licenciatura en Derecho}
\def\documentinstitution{UnADM}
\def\documentmatricula{ES2611202040}
\def\documentgroup{020}
\def\documentcourse{Teoría del Estado y Constitución}
\def\documentteacher{}

\hfwidthwrap{false}
\hfwidthtitle{0.55}
\hfwidthcourse{0.42}

\begin{document}
% Wrapper local para apuntar a la plantilla maestra de AulaTeX
\input{C:/Users/Sysx/Documents/AulaTeX-Academico/aulaTEX/base/Plantilla-Informe/template.tex}


\begin{abstractd}
El presente documento responde al foro de la Semana 8 sobre la razón de la existencia del Estado, analizando los fines más relevantes del Estado moderno, el impacto de la globalización en su estructura y función, y la obligatoriedad de los tratados internacionales en México conforme al artículo 1° constitucional. Se adopta una postura crítica fundamentada en la doctrina jurídica y constitucional mexicana.
\end{abstractd}

\templateIndex

\section{Los fines esenciales del Estado: una perspectiva garantista}

El Estado, como entidad jurídico-política, encuentra su razón de ser en la necesidad de garantizar la convivencia pacífica, la seguridad jurídica y la protección de los derechos fundamentales de las personas. Desde una perspectiva garantista, el fin más relevante del Estado es la protección y promoción de los derechos humanos, entendidos como límites infranqueables al ejercicio del poder público.

La Constitución Política de los Estados Unidos Mexicanos, en su artículo 1°, establece que todas las autoridades tienen la obligación de promover, respetar, proteger y garantizar los derechos humanos. Este mandato constitucional refleja la evolución del Estado mexicano hacia un modelo garantista, donde el respeto a la dignidad humana se erige como el fin último y primordial.

\section{Impacto de la globalización en el Estado moderno}

La globalización ha transformado radicalmente la estructura y funciones del Estado moderno. Este fenómeno ha generado una tensión dialéctica entre la soberanía estatal y la interdependencia global, obligando a los Estados a adaptar sus instituciones y marcos normativos.

En el ámbito jurídico, la globalización ha impulsado la armonización de legislaciones nacionales con estándares internacionales, especialmente en materia de derechos humanos y comercio. Sin embargo, este proceso no está exento de desafíos. La soberanía estatal se ve limitada por la necesidad de cumplir con tratados internacionales y decisiones de organismos supranacionales, lo que genera debates sobre la legitimidad y el alcance de estas restricciones.

En el contexto mexicano, la globalización ha tenido un impacto significativo en la reforma constitucional de derechos humanos de 2011. Esta reforma incorporó el principio pro persona y estableció la obligatoriedad de los tratados internacionales en materia de derechos humanos, elevándolos a rango constitucional.

\section{La obligatoriedad de los tratados internacionales en México}

El artículo 1° constitucional establece que los tratados internacionales en materia de derechos humanos son obligatorios para todas las autoridades mexicanas. Esta disposición refleja un avance significativo en la protección de los derechos fundamentales, al incorporar estándares internacionales al orden jurídico nacional.

Considero que la obligatoriedad de los tratados internacionales es un avance positivo, ya que permite una protección más robusta de los derechos humanos. Sin embargo, es fundamental que esta incorporación se realice bajo un marco de control constitucional que garantice la supremacía de la Constitución y evite conflictos normativos.

La Suprema Corte de Justicia de la Nación ha establecido criterios para resolver posibles contradicciones entre tratados internacionales y la Constitución, privilegiando siempre la protección más favorable a la persona. Este enfoque garantista es esencial para asegurar que la incorporación de tratados internacionales no vulnere principios fundamentales del orden constitucional mexicano.

\section{Conclusiones}
\clearpage

Considero que el fin más relevante del Estado es la protección y promoción de los derechos humanos. Este fin no solo justifica la existencia del Estado, sino que también establece límites claros al ejercicio del poder público, garantizando que todas las acciones estatales se orienten hacia el respeto y la dignidad de las personas.

La globalización ha impactado significativamente al Estado moderno, generando una tensión entre la soberanía estatal y la interdependencia global. Este fenómeno ha impulsado la armonización de legislaciones nacionales con estándares internacionales, especialmente en materia de derechos humanos. Sin embargo, es fundamental que este proceso se realice bajo un marco de control constitucional que garantice la supremacía de la Constitución y evite conflictos normativos.

En cuanto a la obligatoriedad de los tratados internacionales en México, sostengo con firmeza que representa un avance significativo en la protección de los derechos fundamentales. No obstante, es crucial que esta incorporación se realice bajo un marco de control constitucional que garantice la supremacía de la Constitución y evite conflictos normativos. La postura de la Suprema Corte de Justicia de la Nación, que privilegia la protección más favorable a la persona, es un ejemplo de cómo el Estado mexicano puede adaptarse a los desafíos de la globalización sin renunciar a su soberanía constitucional.

\footnote{Declaración: Este trabajo académico fue desarrollado por \documentauthorcanonico\ (matrícula \documentmatricula) como parte de las actividades de la asignatura \documentcourse\ de la \documentsubject\ en la \documentinstitution. Se utilizó apoyo tecnológico para la organización y redacción, pero el contenido, análisis y conclusiones son responsabilidad exclusiva del autor, conforme a los lineamientos institucionales de la UnADM sobre el uso de herramientas de asistencia académica.}

\clearpage
\printbibliography
\end{document}